\section{USAMO 2018}
        \begin{enumerate}
        	\item (\textit{USAMO 2018}) Note: For any geometry problem whose statement begins with an asterisk ($*$), the first page of the solution must be a large, in-scale, clearly labeled diagram. Failure to meet this requirement will result in an automatic 1-point deduction.


 Let $a,b,c$ be positive real numbers such that $a+b+c=4\sqrt[3]{abc}$. Prove that 
 \[2(ab+bc+ca)+4\min(a^2,b^2,c^2)\ge a^2+b^2+c^2.\]


 

	\item (\textit{USAMO 2018}) Note: For any geometry problem whose statement begins with an asterisk ($*$), the first page of the solution must be a large, in-scale, clearly labeled diagram. Failure to meet this requirement will result in an automatic 1-point deduction.


 Find all functions $f:(0,\infty) \rightarrow (0,\infty)$ such that


 
 \[f\left(x+\frac{1}{y}\right)+f\left(y+\frac{1}{z}\right) + f\left(z+\frac{1}{x}\right) = 1\]
for all $x,y,z >0$ with $xyz =1.$


 

	\item (\textit{USAMO 2018}) Note: For any geometry problem whose statement begins with an asterisk ($*$), the first page of the solution must be a large, in-scale, clearly labeled diagram. Failure to meet this requirement will result in an automatic 1-point deduction.


 For a given integer $n\ge 2,$ let $\{a_1,a_2,\ldots,a_m\}$ be the set of positive integers less than $n$ that are relatively prime to $n.$ Prove that if every prime that divides $m$ also divides $n,$ then $a_1^k+a_2^k + \dots + a_m^k$ is divisible by $m$ for every positive integer $k.$


 

	\item (\textit{USAMO 2018}) Note: For any geometry problem whose statement begins with an asterisk ($*$), the first page of the solution must be a large, in-scale, clearly labeled diagram. Failure to meet this requirement will result in an automatic 1-point deduction.


 Let $p$ be a prime, and let $a_1, \dots, a_p$ be integers. Show that there exists an integer $k$ such that the numbers 
 \[a_1 + k, a_2 + 2k, \dots, a_p + pk\]produce at least $\frac{1}{2} p$ distinct remainders upon division by $p$.


 

	\item (\textit{USAMO 2018}) Note: For any geometry problem whose statement begins with an asterisk ($*$), the first page of the solution must be a large, in-scale, clearly labeled diagram. Failure to meet this requirement will result in an automatic 1-point deduction.


 In convex cyclic quadrilateral $ABCD,$ we know that lines $AC$ and $BD$ intersect at $E,$ lines $AB$ and $CD$ intersect at $F,$ and lines $BC$ and $DA$ intersect at $G.$ Suppose that the circumcircle of $\triangle ABE$ intersects line $CB$ at $B$ and $P$, and the circumcircle of $\triangle ADE$ intersects line $CD$ at $D$ and $Q$, where $C,B,P,G$ and $C,Q,D,F$ are collinear in that order. Prove that if lines $FP$ and $GQ$ intersect at $M$, then $\angle MAC = 90^{\circ}.$


 

	\item (\textit{USAMO 2018}) Note: For any geometry problem whose statement begins with an asterisk ($*$), the first page of the solution must be a large, in-scale, clearly labeled diagram. Failure to meet this requirement will result in an automatic 1-point deduction.


 Let $a_n$ be the number of permutations $(x_1, x_2, \dots, x_n)$ of the numbers $(1,2,\dots, n)$ such that the $n$ ratios $\frac{x_k}{k}$ for $1\le k\le n$ are all distinct. Prove that $a_n$ is odd for all $n\ge 1.$


 

\end{enumerate}